%----------------------------------------------------------------------------------------
% GONIOMETRISCHE VERGELIJKINGEN. Alle opgaven zijn eigen werk.
%
% Opbouw: Start, leerstof (theorie, werkwijze of uitleg, voorbeeld, korte oefening), en op een
% nieuwe pagina de Opdrachten (door elkaar, stijgende moeilijkheid), Koppel en Test jezelf.
% Labels van dit hoofdstuk beginnen met h4-.
%----------------------------------------------------------------------------------------

\mychapter{4}{Goniometrische vergelijkingen}

\lkr{Leerplan Basis: LPD 11-13 en K2 (goniometrische functies), het deel over de cyclometrische functies en goniometrische vergelijkingen. De algemene oplossing met $k \in \Z$ is voor veel leerlingen abstract: laat ze telkens ook de oplossingen in een interval opschrijven, en controleren op de eenheidscirkel. Het rekentoestel geeft maar één oplossing; de tweede vinden ze met de verwante hoeken.}

Wanneer is het water in de haven diep genoeg? Wanneer is de gondel hoger dan 100 m? Zulke vragen leiden tot een vergelijking waarin de onbekende in een sinus of cosinus staat. Omdat de functies periodiek zijn, hebben zulke vergelijkingen meestal oneindig veel oplossingen.

\begin{oefeningen}[Start]
\oef{h4-start-1}

\oef{h4-verken}
\end{oefeningen}

%========================================================================================
\section{Cyclometrische functies}

De sinus is niet injectief: $\sin\frac{\pi}{6} = \sin\frac{5\pi}{6}$. Beperk je het domein tot $\bigl[-\frac{\pi}{2}, \frac{\pi}{2}\bigr]$, waar de sinus stijgt van $-1$ tot $1$, dan heeft ze wel een inverse (zie Algemene functies). Net zo voor cosinus en tangens.

\begin{theorie}
\renewcommand{\arraystretch}{1.35}
\begin{tabular}{@{}l@{\hspace{5mm}}l@{\hspace{5mm}}l@{\hspace{5mm}}l@{}}
functie & inverse van & domein & bereik \\ \hline
$\formule{\bgsin x}$ & $\sin$ op $\bigl[-\frac{\pi}{2}, \frac{\pi}{2}\bigr]$ & $[-1, 1]$ & $\formule{\bigl[-\frac{\pi}{2}, \frac{\pi}{2}\bigr]}$ \\
$\formule{\bgcos x}$ & $\cos$ op $[0, \pi]$ & $[-1, 1]$ & $\formule{[0, \pi]}$ \\
$\formule{\bgtan x}$ & $\tan$ op $\bigl]-\frac{\pi}{2}, \frac{\pi}{2}\bigr[$ & $\R$ & $\formule{\bigl]-\frac{\pi}{2}, \frac{\pi}{2}\bigr[}$
\end{tabular}

\vspace{1mm}
Dus $\formule{\bgsin a = \theta \iff \sin\theta = a \text{ en } -\tfrac{\pi}{2} \leqslant \theta \leqslant \tfrac{\pi}{2}}$, en zo verder.
\end{theorie}

\naast{%
De naam komt van \textit{boogsinus}: $\bgsin a$ is de boog (de hoek in radialen) waarvan de sinus $a$ is. Op het rekentoestel heet de toets $\sin^{-1}$; zet het toestel in radialen.

De grafiek van $\bgsin$ is het spiegelbeeld van het stijgende stuk van de sinus om de rechte $y = x$.}{%
\begin{tikzpicture}[x=.85cm,y=.85cm,font=\scriptsize]
  \draw[->,grijs] (-1.9,0) -- (2,0) node[right,text=black] {$x$};
  \draw[->,grijs] (0,-1.9) -- (0,2) node[above,text=black] {$y$};
  \draw[grijs,dashed] (-1.8,-1.8) -- (1.8,1.8);
  \draw[accent,thick,dashed] plot[domain=-1.5708:1.5708,samples=50,smooth] (\x,{sin(\x r)});
  \draw[accent,thick] plot[domain=-1:1,samples=80,smooth] (\x,{asin(\x)*pi/180});
  \node[accent,right] at (1,1.5708) {$\bgsin x$};
  \node[accent,below right] at (1.4,.9) {$\sin x$};
  \node[left] at (0,1.5708) {$\frac{\pi}{2}$};
\end{tikzpicture}}

\voorbeeld{}{Bepaal exact $\bgcos\bigl(-\frac{\sqrt2}{2}\bigr)$.\label{vb:h4-bg}}{%
\uitleg{Zoek de hoek in $[0, \pi]$ met cosinus $-\frac{\sqrt2}{2}$. In dat interval is de cosinus negatief in het tweede kwadrant.}
\begin{stappenbreed}
\cos\theta = -\tfrac{\sqrt2}{2},\ 0 \leqslant \theta \leqslant \pi & Definitie. \\
\theta = \pi - \tfrac{\pi}{4} = \tfrac{3\pi}{4} & Supplementair met $\frac{\pi}{4}$.
\end{stappenbreed}}

\nujij{Bepaal exact $\bgsin\bigl(-\frac12\bigr)$ en $\bgtan 1$.}{$-\frac{\pi}{6}$, want $\sin\bigl(-\frac{\pi}{6}\bigr) = -\frac12$ en $-\frac{\pi}{6}$ ligt in $\bigl[-\frac{\pi}{2}, \frac{\pi}{2}\bigr]$. $\bgtan 1 = \frac{\pi}{4}$.}

\begin{oefeningen}
\oef{h4-bg}
\end{oefeningen}

%========================================================================================
\section{Elementaire vergelijkingen}

\naast{%
Een horizontale rechte $y = a$ (met $-1 < a < 1$) snijdt de eenheidscirkel in twee punten: bij de hoek $\bgsin a$ en bij de supplementaire hoek $\pi - \bgsin a$. Daarbij komen telkens veelvouden van $2\pi$. Voor de cosinus gebruik je een verticale rechte: de twee hoeken zijn tegengesteld.}{%
\begin{tikzpicture}[scale=1.3,font=\scriptsize]
  \eenheidscirkel
  \draw[accent] (-1.2,.45) -- (1.25,.45) node[right] {$y = a$};
  \draw[accent] (0,0) -- (26.7:1) (0,0) -- (153.3:1);
  \fill[accent] (26.7:1) circle (.035);
  \fill[accent] (153.3:1) circle (.035);
  \node[accent] at (12:.55) {$\theta$};
  \node[accent,above] at (153.3:1.05) {$\pi - \theta$};
\end{tikzpicture}}

\begin{theorie}
\begin{stelling}{Elementaire vergelijkingen}
Met $k \in \Z$:
\[\formule{\sin x = a \iff x = \bgsin a + k \cdot 2\pi \ \text{ of } \ x = \pi - \bgsin a + k \cdot 2\pi} \qquad (-1 \leqslant a \leqslant 1)\]
\[\formule{\cos x = a \iff x = \bgcos a + k \cdot 2\pi \ \text{ of } \ x = -\bgcos a + k \cdot 2\pi} \qquad (-1 \leqslant a \leqslant 1)\]
\[\formule{\tan x = a \iff x = \bgtan a + k\pi}\]
Is $|a| > 1$, dan hebben $\sin x = a$ en $\cos x = a$ geen oplossingen.
\end{stelling}
\end{theorie}

\voorbeeld{}{Los op: $\sin x = \frac12$. Geef de algemene oplossing en de oplossingen in $[0, 2\pi[$.\label{vb:h4-sin}}{%
\uitleg{$\bgsin\frac12 = \frac{\pi}{6}$; de tweede oplossing is de supplementaire hoek.}
\begin{stappenbreed}
x = \tfrac{\pi}{6} + k \cdot 2\pi \ \text{ of } \ x = \tfrac{5\pi}{6} + k \cdot 2\pi & $\pi - \frac{\pi}{6} = \frac{5\pi}{6}$, $k \in \Z$. \\
x \in \bigl\{\tfrac{\pi}{6}, \tfrac{5\pi}{6}\bigr\} & In $[0, 2\pi[$: $k = 0$.
\end{stappenbreed}}

Staat er een veelvoud van $x$ in de sinus of cosinus, los dan eerst op naar dat veelvoud, en deel daarna, \textit{ook de periode}.

\voorbeeld{}{Los op: $\cos 2x = -\frac12$. Geef de oplossingen in $[0, 2\pi[$.\label{vb:h4-cos}}{%
\uitleg{Los op naar $2x$, en deel dan door 2. De periode $2\pi$ wordt $\pi$: er zijn dus twee keer zoveel oplossingen in $[0, 2\pi[$.}
\begin{stappenbreed}
2x = \tfrac{2\pi}{3} + k \cdot 2\pi \ \text{ of } \ 2x = -\tfrac{2\pi}{3} + k \cdot 2\pi & $\bgcos\bigl(-\frac12\bigr) = \frac{2\pi}{3}$. \\
x = \tfrac{\pi}{3} + k\pi \ \text{ of } \ x = -\tfrac{\pi}{3} + k\pi & Delen door 2. \\
x \in \bigl\{\tfrac{\pi}{3}, \tfrac{2\pi}{3}, \tfrac{4\pi}{3}, \tfrac{5\pi}{3}\bigr\} & $k = 0, 1$ en $k = 1, 2$.
\end{stappenbreed}}

\begin{werkwijze}
Zonder $\sin$, $\cos$ of $\tan$ af. Zoek één oplossing met de tabel of met $\bgsin$, $\bgcos$, $\bgtan$ (rekentoestel in radialen). Schrijf de tweede oplossing met de verwante hoek en tel $k \cdot 2\pi$ op (bij de tangens: $k\pi$). Los daarna op naar $x$. Zoek de oplossingen in het gevraagde interval door $k = 0, 1, -1, \ldots$ in te vullen.
\end{werkwijze}

\begin{oefeningen}
\oef{h4-elementair}
\end{oefeningen}

%========================================================================================
\section{Herleidbare vergelijkingen}

Andere vergelijkingen breng je met de formules uit hoofdstuk 1 en 2 terug tot elementaire vergelijkingen.

\voorbeeld{}{Los op: $2\sin^2 x - \sin x - 1 = 0$.\label{vb:h4-kwadratisch}}{%
\uitleg{Een tweedegraadsvergelijking in $\sin x$. Stel $t = \sin x$.}
\begin{stappenbreed}
2t^2 - t - 1 = 0 \iff t = 1 \text{ of } t = -\tfrac12 & Discriminant $9$. \\
\sin x = 1 \iff x = \tfrac{\pi}{2} + k \cdot 2\pi & \\
\sin x = -\tfrac12 \iff x = -\tfrac{\pi}{6} + k \cdot 2\pi \text{ of } x = \tfrac{7\pi}{6} + k \cdot 2\pi & $\pi - \bigl(-\frac{\pi}{6}\bigr) = \frac{7\pi}{6}$.
\end{stappenbreed}}

\voorbeeld{}{Los op: $\sin 2x = \cos x$.\label{vb:h4-ontbinden}}{%
\uitleg{Schrijf $\sin 2x$ met de verdubbelingsformule, breng alles naar één lid en ontbind. Deel niet door $\cos x$: dan verlies je oplossingen.}
\begin{stappenbreed}
2\sin x\cos x - \cos x = 0 \iff \cos x\,(2\sin x - 1) = 0 & Gemeenschappelijke factor. \\
\cos x = 0 \iff x = \tfrac{\pi}{2} + k\pi & \\
\sin x = \tfrac12 \iff x = \tfrac{\pi}{6} + k \cdot 2\pi \text{ of } x = \tfrac{5\pi}{6} + k \cdot 2\pi &
\end{stappenbreed}}

\voorbeeld{}{Los op: $\sin x + \sqrt3\cos x = 1$.\label{vb:h4-samenvoegen}}{%
\uitleg{Voeg samen tot één sinus (hoofdstuk 2): $\sin x + \sqrt3\cos x = 2\sin\bigl(x + \frac{\pi}{3}\bigr)$.}
\begin{stappenbreed}
\sin\bigl(x + \tfrac{\pi}{3}\bigr) = \tfrac12 & \\
x + \tfrac{\pi}{3} = \tfrac{\pi}{6} + k \cdot 2\pi \ \text{ of } \ x + \tfrac{\pi}{3} = \tfrac{5\pi}{6} + k \cdot 2\pi & \\
x = -\tfrac{\pi}{6} + k \cdot 2\pi \ \text{ of } \ x = \tfrac{\pi}{2} + k \cdot 2\pi & Controle: $\sin\frac{\pi}{2} + \sqrt3\cos\frac{\pi}{2} = 1$.
\end{stappenbreed}}

\begin{oefeningen}
\oef{h4-herleidbaar}
\end{oefeningen}

%========================================================================================
\section{Vergelijkingen in een context}

\voorbeeld{}{In de haven uit hoofdstuk 3 is de waterhoogte $h(t) = 2{,}2\cos\bigl(\frac{2\pi}{12{,}5}(t - 2)\bigr) + 2{,}6$ (in m, $t$ in uur na middernacht). Een schip heeft minstens $3$ m water nodig. Wanneer kan het tussen middernacht en 18 uur binnenvaren?\label{vb:h4-context}}{%
\uitleg{Los eerst de vergelijking $h(t) = 3$ op. Rond de top (om 2 uur) is het water hoger dan $3$ m; de grenzen liggen symmetrisch rond die top.}
\begin{stappenbreed}
\cos\bigl(\tfrac{2\pi}{12{,}5}(t - 2)\bigr) = \tfrac{0{,}4}{2{,}2} \approx 0{,}1818 & \\
\tfrac{2\pi}{12{,}5}(t - 2) = \pm 1{,}388 & $\bgcos 0{,}1818 \approx 1{,}388$, met $k = 0$. \\
t - 2 = \pm 2{,}76 & Vermenigvuldigen met $\frac{12{,}5}{2\pi}$. \\
t \in [0;\ 4{,}76] \text{ en } [11{,}74;\ 17{,}26] & Tel er een periode $12{,}5$ bij voor de tweede.
\end{stappenbreed}
Het schip kan binnenvaren tot ongeveer 4.45 uur, en opnieuw tussen ongeveer 11.45 en 17.15 uur.}

\begin{oefeningen}
\oef{h4-context}
\end{oefeningen}

%========================================================================================
\opdrachtenkop

\begin{oefeningen}
\oef{h4-o1}

\oef{h4-o2}

\oef{h4-o3}

\oef{h4-o4}

\oef{h4-o5}

\oef{h4-o6}

\oef{h4-o7}

\oef{h4-o8}

\oef{h4-o9}
\end{oefeningen}

%----------------------------------------------------------------------------------------
% KOPPEL: verbanden met andere leerstof
%----------------------------------------------------------------------------------------
\needspace{20\baselineskip}
\opdrachtsectie{Koppel}

\lkr{Algemene functies: kijk na wat al gezien is. Het pakket legt de volgorde niet vast.}

\begin{oefeningen}

\oef{h4-koppel-2}

\end{oefeningen}

%----------------------------------------------------------------------------------------
\opdrachtsectie{Test jezelf}

\begin{oefeningen}
\oef{h4-t1}

\oef{h4-t2}

\oef{h4-t3}

\oef{h4-t4}

\oef{h4-t5}
\end{oefeningen}

\herhaalschema{%
  \fv{oef:h4-t1}{vb:h4-bg}{oef:h4-bg}%
  \fv{oef:h4-t2}{vb:h4-sin}{oef:h4-elementair}%
  \fv{oef:h4-t3}{vb:h4-cos}{oef:h4-elementair}%
  \fv{oef:h4-t4}{vb:h4-ontbinden}{oef:h4-herleidbaar}%
  \fv{oef:h4-t5}{vb:h4-context}{oef:h4-context}}

\ifzelfstudie
\section*{Maak je eigen samenvatting}
Sluit dit boek en schrijf in je schrift: domein en bereik van $\bgsin$, $\bgcos$ en $\bgtan$, de algemene oplossing van $\sin x = a$, $\cos x = a$ en $\tan x = a$ met een figuur van de eenheidscirkel, hoe je oplossingen in een interval vindt, en de drie manieren om een vergelijking te herleiden (substitutie, ontbinden, samenvoegen). Vergelijk daarna met de kaders.
\fi
